% ============================================================
%  BLE IMU Guide - XIAO nRF52840 Sense
%  Guia exhaustiva desde cero absoluto
% ============================================================
\documentclass[11pt,a4paper]{article}

\usepackage[utf8]{inputenc}
\usepackage[T1]{fontenc}
\usepackage[spanish,es-noquoting,es-noshorthands]{babel}
\usepackage{geometry}
\geometry{margin=2.5cm}
\usepackage{parskip}
\usepackage{enumitem}
\usepackage{hyperref}
\hypersetup{colorlinks=true,linkcolor=blue,urlcolor=blue}
\usepackage{listings}
\usepackage{xcolor}
\usepackage{graphicx}
\usepackage{amsmath}
\usepackage{fancyhdr}
\usepackage{titlesec}

% ── Estilos de codigo ────────────────────────────────────
\definecolor{codebg}{HTML}{F5F5F5}
\definecolor{codegreen}{HTML}{2E7D32}
\definecolor{codegray}{HTML}{757575}
\definecolor{codeblue}{HTML}{1565C0}

\lstdefinestyle{arduino}{
    language=C++,
    backgroundcolor=\color{codebg},
    basicstyle=\ttfamily\small,
    breaklines=true,
    frame=single,
    rulecolor=\color{gray!40},
    keywordstyle=\color{codeblue}\bfseries,
    stringstyle=\color{codegreen},
    commentstyle=\color{codegray}\itshape,
    numberstyle=\tiny\color{codegray},
    numbers=left,
    numbersep=8pt,
    tabsize=2,
    showstringspaces=false,
    morekeywords={BLEService,BLECharacteristic,BLEDevice,LSM6DS3,
                  BLERead,BLENotify,uint8_t,float,void,byte},
}

\lstdefinestyle{python}{
    language=Python,
    backgroundcolor=\color{codebg},
    basicstyle=\ttfamily\small,
    breaklines=true,
    frame=single,
    rulecolor=\color{gray!40},
    keywordstyle=\color{codeblue}\bfseries,
    stringstyle=\color{codegreen},
    commentstyle=\color{codegray}\itshape,
    numberstyle=\tiny\color{codegray},
    numbers=left,
    numbersep=8pt,
    tabsize=4,
    showstringspaces=false,
    morekeywords={async,await,asyncio,struct,bytearray,BleakClient,
                  BleakScanner,True,False,None,f},
}

\pagestyle{fancy}
\fancyhf{}
\fancyhead[L]{\small\textit{SensorKit --- BLE IMU Guide}}
\fancyhead[R]{\small\thepage}
\renewcommand{\headrulewidth}{0.4pt}

\title{%
    \vspace{-1cm}
    \textbf{Guia Completa: Lectura de IMU via BLE}\\[6pt]
    \large Seeed XIAO nRF52840 Sense $\rightarrow$ Python\\[4pt]
    \normalsize Todo explicado desde cero, sin dar nada por sabido
}
\author{SensorKit Project}
\date{Marzo 2026}

\begin{document}
\maketitle
\tableofcontents
\newpage


% ============================================================
\part{Fundamentos: todo lo que necesitas saber antes de tocar codigo}
% ============================================================


% ============================================================
\section{Que es lo que queremos hacer (y por que es complicado)}
% ============================================================

Nuestro objetivo final es muy simple de describir: queremos leer
numeros de un sensor y verlos en nuestro ordenador. Pero hay un
problema: el sensor no esta enchufado al ordenador con un cable.
Esta en una plaquita pequeña (la XIAO nRF52840 Sense) que se
comunica por Bluetooth.

Esto introduce varias capas de complejidad:

\begin{enumerate}
    \item El sensor no ``habla'' Bluetooth. El sensor habla un
          protocolo llamado I$^2$C (lo veremos despues). Asi que
          necesitamos un \textbf{traductor}: el microcontrolador
          nRF52840 que esta en la placa, que lee los datos del
          sensor por I$^2$C y los reempaqueta para enviarlos por
          Bluetooth.

    \item Bluetooth no es como un cable USB donde simplemente
          ``fluyen'' bytes. Bluetooth Low Energy (BLE) tiene una
          estructura rigida: los datos se organizan en ``servicios''
          y ``caracteristicas'' (esto se llama GATT, y lo
          explicaremos a fondo).

    \item Del lado del ordenador, necesitamos un programa que sepa
          hablar BLE, conectarse al dispositivo, y extraer los
          datos. Para esto usamos Python con la libreria
          \texttt{bleak}.
\end{enumerate}

El flujo completo es:

\begin{verbatim}
  SENSOR (chip LSM6DS3TR-C)
    "Tengo una aceleracion de 0.02g en X"
         |
         | Protocolo I2C (cables fisicos dentro de la placa)
         v
  MICROCONTROLADOR (chip nRF52840)
    "Empaqueto esos datos como bytes y los pongo
     disponibles via Bluetooth"
         |
         | Ondas de radio Bluetooth (2.4 GHz, sin cables)
         v
  TU ORDENADOR (Mac/PC/Linux)
    "Me conecto por Bluetooth, leo los bytes,
     los convierto de vuelta a numeros"
         |
         | Libreria bleak + struct de Python
         v
  TU SCRIPT PYTHON
    "Ahora tengo ax=0.02, ay=0.01, az=1.00
     y puedo hacer lo que quiera con ellos"
\end{verbatim}

Cada una de estas flechas es un mundo. Vamos a explicarlas todas.


% ============================================================
\section{El sensor: LSM6DS3TR-C}
% ============================================================

\subsection{Que es un sensor IMU}

IMU significa \textbf{Inertial Measurement Unit} (Unidad de Medicion
Inercial). Es un chip que mide movimiento. El LSM6DS3TR-C, fabricado
por la empresa \textbf{STMicroelectronics}, tiene dos sensores dentro:

\begin{itemize}
    \item \textbf{Acelerometro}: Mide \emph{aceleracion lineal}. Si
          mueves la placa hacia la derecha, el acelerometro detecta
          esa aceleracion. Tambien detecta la gravedad terrestre: si
          la placa esta quieta sobre una mesa, el eje que apunta
          hacia arriba mide aproximadamente $1g$ (donde
          $1g = 9.81 \text{ m/s}^2$). Los otros dos ejes miden
          aproximadamente $0g$.

          El acelerometro funciona internamente con una masa
          microscopica suspendida por muelles de silicio. Cuando
          el chip acelera, la masa se queda atras por inercia y
          deforma los muelles. La deformacion se mide como un
          cambio en la capacitancia electrica entre placas
          microscopicas. Ese cambio de capacitancia se convierte
          en un voltaje, y luego un conversor analogico-digital
          (ADC) lo convierte en un numero.

    \item \textbf{Giroscopio}: Mide \emph{velocidad angular}
          (que tan rapido estas girando). Si giras la placa sobre
          su eje Z, el giroscopio del eje Z mide esa velocidad
          en \textbf{grados por segundo} (dps). Si la placa esta
          quieta, todos los ejes miden aproximadamente 0.

          El giroscopio usa el efecto Coriolis: dentro del chip hay
          una estructura que vibra a alta frecuencia. Cuando el chip
          gira, la fuerza de Coriolis desvía la vibracion en una
          direccion perpendicular. Esa desviacion se mide como un
          cambio de capacitancia, se convierte en voltaje, y luego
          en un numero.
\end{itemize}

\subsection{Ejes X, Y, Z}

Cada sensor mide en \textbf{tres ejes perpendiculares}: X, Y, Z.
Imagina que pones la placa sobre la mesa:

\begin{itemize}
    \item \textbf{X}: Apunta hacia la derecha de la placa.
    \item \textbf{Y}: Apunta hacia ``arriba'' de la placa (el lado
          opuesto al conector USB).
    \item \textbf{Z}: Apunta hacia arriba, perpendicular a la mesa.
\end{itemize}

En total tenemos 6 valores: $a_x, a_y, a_z$ (acelerometro) y
$g_x, g_y, g_z$ (giroscopio). Por eso se dice ``6 ejes'' o ``6-DOF''
(6 Degrees of Freedom).

\subsection{Rangos y precision}

El sensor tiene rangos configurables. Un rango mayor permite medir
movimientos mas bruscos, pero con menos precision:

\begin{center}
\begin{tabular}{lll}
\hline
\textbf{Sensor} & \textbf{Rangos disponibles} & \textbf{Default} \\
\hline
Acelerometro & $\pm$2g, $\pm$4g, $\pm$8g, $\pm$16g & $\pm$2g \\
Giroscopio & $\pm$125, $\pm$245, $\pm$500, $\pm$1000, $\pm$2000 dps & $\pm$245 dps \\
\hline
\end{tabular}
\end{center}

Internamente, el sensor convierte su medicion a un numero entero de
16 bits (valores de $-32768$ a $+32767$). Si el rango del
acelerometro es $\pm$2g, entonces $+32767$ corresponde a $+2g$ y
$-32768$ corresponde a $-2g$. La resolucion es:
$$\frac{2g}{32768} \approx 0.000061 \text{ g por unidad}$$

La libreria que usamos en Arduino hace esta conversion por nosotros
y nos da directamente el valor en $g$ o $dps$.

\subsection{Registros: como se comunica el sensor}

El sensor LSM6DS3TR-C es, en esencia, una pequeña memoria con
``casillas'' llamadas \textbf{registros}. Cada registro tiene una
direccion (un numero hexadecimal) y almacena un byte (8 bits,
valor de 0 a 255).

Hay dos tipos de registros:

\begin{itemize}
    \item \textbf{Registros de configuracion}: Tu escribes en ellos
          para decirle al sensor como comportarse. Por ejemplo, el
          registro \texttt{CTRL1\_XL} (direccion \texttt{0x10})
          controla la frecuencia de muestreo y el rango del
          acelerometro. Si escribes el byte \texttt{0x60} en ese
          registro, le estas diciendo: ``muestrea a 416 Hz con
          rango $\pm$2g''.

    \item \textbf{Registros de datos}: El sensor escribe en ellos
          los valores medidos. Tu los lees. Los datos del
          acelerometro eje X estan en los registros \texttt{0x28}
          (byte bajo) y \texttt{0x29} (byte alto). Juntos forman
          un entero de 16 bits.
\end{itemize}

\begin{center}
\begin{tabular}{lll}
\hline
\textbf{Registro} & \textbf{Direccion} & \textbf{Contenido} \\
\hline
WHO\_AM\_I & 0x0F & Siempre vale 0x6A (identifica al chip) \\
CTRL1\_XL & 0x10 & Configuracion del acelerometro \\
CTRL2\_G & 0x11 & Configuracion del giroscopio \\
OUTX\_L\_XL & 0x28 & Acelerometro X, byte bajo \\
OUTX\_H\_XL & 0x29 & Acelerometro X, byte alto \\
OUTY\_L\_XL & 0x2A & Acelerometro Y, byte bajo \\
... & ... & ... \\
OUTX\_L\_G & 0x22 & Giroscopio X, byte bajo \\
OUTX\_H\_G & 0x23 & Giroscopio X, byte alto \\
... & ... & ... \\
\hline
\end{tabular}
\end{center}

Pero nosotros no leemos estos registros directamente en nuestro
codigo. La libreria \texttt{Seeed\_Arduino\_LSM6DS3} lo hace por
nosotros. Cuando llamamos a \texttt{myIMU.readFloatAccelX()}, la
libreria:
\begin{enumerate}
    \item Lee los registros \texttt{0x28} y \texttt{0x29}.
    \item Combina los dos bytes en un entero de 16 bits.
    \item Multiplica por el factor de escala ($0.000061$ para $\pm$2g).
    \item Devuelve el resultado como un \texttt{float} en unidades de $g$.
\end{enumerate}


% ============================================================
\section{I$^2$C: como el microcontrolador habla con el sensor}
% ============================================================

\subsection{Que es I$^2$C}

I$^2$C (pronunciado ``i-cuadrado-ce'' o ``i-dos-ce'') significa
\textbf{Inter-Integrated Circuit}. Es un protocolo de comunicacion
que permite a chips electronicos hablar entre si usando solo
\textbf{dos cables}:

\begin{itemize}
    \item \textbf{SDA} (Serial Data): Por aqui viajan los datos
          en ambas direcciones (el micro le envia comandos al sensor,
          el sensor le devuelve los datos al micro).
    \item \textbf{SCL} (Serial Clock): Es una senal de reloj que
          marca el ritmo. Cada ``tick'' del reloj significa
          ``ahora lee un bit de SDA''. Esto sincroniza a ambos
          chips para que no se pierdan datos.
\end{itemize}

\subsection{Master y Slave}

En I$^2$C siempre hay un \textbf{master} (quien manda) y uno o
mas \textbf{slaves} (quienes obedecen):

\begin{itemize}
    \item \textbf{Master}: El microcontrolador nRF52840. El inicia
          todas las comunicaciones. Dice ``quiero leer el registro
          0x28 del dispositivo 0x6A''.
    \item \textbf{Slave}: El sensor LSM6DS3TR-C. Su direccion en
          el bus es \textbf{0x6A}. Cuando el master pide el
          registro 0x28, el sensor responde con el byte que tiene
          almacenado ahi.
\end{itemize}

\subsection{Direccion 0x6A}

Cada dispositivo I$^2$C tiene una direccion unica de 7 bits (valores
de 0 a 127, o 0x00 a 0x7F en hexadecimal). Esto permite tener
hasta 128 dispositivos en el mismo bus (mismos dos cables). El
master incluye la direccion en cada comunicacion para que solo el
dispositivo correcto responda.

El LSM6DS3TR-C tiene la direccion \textbf{0x6A} porque asi lo
diseño STMicroelectronics. Viene fija de fabrica (bueno, se puede
cambiar a 0x6B moviendo un pin, pero en la XIAO esta soldado a
0x6A).

\subsection{Como se ve una lectura I$^2$C por dentro}

Cuando el micro quiere leer el eje X del acelerometro, la
conversacion en el bus I$^2$C es algo asi:

\begin{enumerate}
    \item El master envia una condicion de START (baja SDA mientras
          SCL esta alto --- esta es una senal electrica especial
          que todos los slaves reconocen como ``atencion, empieza
          una comunicacion'').
    \item El master envia la direccion 0x6A + bit de escritura.
          El slave responde con ACK (``estoy aqui'').
    \item El master envia 0x28 (``quiero empezar a leer desde el
          registro 0x28'').
    \item El master envia otra condicion de START (repeated start).
    \item El master envia la direccion 0x6A + bit de lectura.
    \item El slave envia el contenido del registro 0x28 (byte bajo
          del acelerometro X).
    \item El master envia ACK (``dame mas'').
    \item El slave envia el contenido de 0x29 (byte alto del
          acelerometro X).
    \item El master envia NACK (``ya no quiero mas'') y condicion
          de STOP.
\end{enumerate}

Todo esto ocurre a una velocidad de 400 kHz (400,000 ciclos de
reloj por segundo), asi que una lectura completa tarda
microsegundos. Nosotros no vemos nada de esto en nuestro codigo
porque la libreria \texttt{Wire} de Arduino lo hace todo por
nosotros.

En la XIAO, los cables SDA y SCL entre el nRF52840 y el LSM6DS3TR-C
estan \textbf{dentro de la placa}, son pistas de cobre en el
circuito impreso. No tienes que conectar nada.


% ============================================================
\section{Bluetooth Low Energy (BLE): como los datos viajan al ordenador}
% ============================================================

\subsection{Que es Bluetooth y por que ``Low Energy''}

Bluetooth es un estandar de comunicacion inalambrica por radio,
operando en la banda de \textbf{2.4 GHz} (la misma frecuencia que
el WiFi). Fue inventado en los años 90 para reemplazar cables cortos
(auriculares, teclados, etc.).

Existen dos ``sabores'' de Bluetooth:

\begin{itemize}
    \item \textbf{Bluetooth Clasico (BR/EDR)}: Diseñado para
          transferir mucha informacion continua. Se usa para audio
          (auriculares), transferencia de archivos, etc. Consume
          bastante energia.

    \item \textbf{Bluetooth Low Energy (BLE)}: Introducido en
          Bluetooth 4.0 (2010). Diseñado para dispositivos que
          envian \emph{poquitos datos de vez en cuando} y necesitan
          que la bateria dure meses o años. Se usa en pulseras de
          actividad, sensores, beacons, etc. \textbf{Es lo que
          usamos nosotros.}
\end{itemize}

BLE no es una ``version menor'' del Bluetooth clasico. Es un
protocolo \textbf{completamente diferente}, incompatible con el
clasico. Un dispositivo BLE no puede hablar con uno clasico y
viceversa (aunque muchos chips modernos, como el nRF52840,
soportan ambos).

\subsection{Roles en BLE: Peripheral y Central}

En BLE siempre hay dos roles:

\begin{itemize}
    \item \textbf{Peripheral (periferico)}: Es el dispositivo
          ``pequeño'' que \emph{ofrece datos}. En nuestro caso,
          la XIAO. El periferico hace dos cosas:

    \begin{enumerate}
        \item \textbf{Advertising}: Antes de que nadie se conecte,
              el periferico grita al aire ``!Estoy aqui! Me llamo
              XIAO-IMU y tengo datos de sensores!'' repetidamente.
              Esto lo hace emitiendo pequeños paquetes de radio
              (31 bytes maximo) cada cierto intervalo (tipicamente
              cada 100 milisegundos). Cualquiera que este escuchando
              puede ver estos anuncios.

        \item \textbf{Servidor GATT}: Una vez que alguien se conecta,
              el periferico actua como un ``servidor'' que tiene
              datos organizados en una estructura (la tabla GATT)
              para que el central pueda leerlos.
    \end{enumerate}

    \item \textbf{Central}: Es el dispositivo ``grande'' que
          \emph{consume datos}. En nuestro caso, tu ordenador.
          El central:

    \begin{enumerate}
        \item \textbf{Escanea}: Escucha los paquetes de advertising
              para descubrir que perifericos hay cerca.
        \item \textbf{Se conecta}: Elige un periferico y establece
              una conexion.
        \item \textbf{Lee datos}: Accede a la tabla GATT del
              periferico para obtener los valores de los sensores.
    \end{enumerate}
\end{itemize}

\textbf{Analogia}: Piensa en una tienda. El periferico es la tienda
que pone un cartel en la puerta (advertising) diciendo ``Tenemos
frutas''. El central es el cliente que ve el cartel, entra (conexion),
y mira el catalogo organizado (GATT) para pedir lo que quiere.


% ============================================================
\section{GATT: la estructura que organiza los datos en BLE}
% ============================================================

\subsection{Que es GATT y por que existe}

GATT significa \textbf{Generic Attribute Profile}. Es la forma en
que BLE organiza los datos que un dispositivo ofrece.

\textbf{El problema que resuelve GATT:} Imagina que tu sensor envia
bytes crudos por Bluetooth. Tu ordenador recibe, por ejemplo, 24
bytes. ¿Que significan? ¿Los primeros 12 son el acelerometro y los
siguientes 12 el giroscopio? ¿O son 6 valores de 4 bytes cada uno?
¿O es todo un solo string de texto? Sin estructura, es imposible
saberlo.

GATT resuelve esto creando una \textbf{jerarquia organizada} donde
cada dato tiene:
\begin{itemize}
    \item Un identificador unico (UUID)
    \item Un tipo y tamaño definido
    \item Propiedades que dicen que puedes hacer con el (leerlo,
          escribirlo, recibir actualizaciones, etc.)
\end{itemize}

\subsection{La jerarquia GATT: Perfil, Servicio, Caracteristica}

La estructura tiene tres niveles, de mas general a mas especifico:

\begin{verbatim}
  SERVIDOR GATT (todo el dispositivo)
  |
  +-- SERVICIO "IMU" (UUID: 19B10000-...)
  |   |
  |   +-- CARACTERISTICA "Acelerometro"
  |   |     UUID: 19B10001-...
  |   |     Valor: 12 bytes
  |   |     Propiedades: READ, NOTIFY
  |   |     +-- DESCRIPTOR CCCD (UUID: 0x2902)
  |   |           Valor: 0x0000 o 0x0001
  |   |
  |   +-- CARACTERISTICA "Giroscopio"
  |         UUID: 19B10002-...
  |         Valor: 12 bytes
  |         Propiedades: READ, NOTIFY
  |         +-- DESCRIPTOR CCCD (UUID: 0x2902)
  |               Valor: 0x0000 o 0x0001
  |
  +-- SERVICIO "Generic Access" (estandar BLE)
  |     (nombre del dispositivo, apariencia, etc.)
  |
  +-- SERVICIO "Generic Attribute" (estandar BLE)
        (control interno de GATT)
\end{verbatim}

Vamos a explicar cada nivel con calma:

\subsubsection{Servidor GATT}

Es todo el dispositivo. La XIAO entera actua como un servidor GATT.
Cuando tu ordenador se conecta, lo primero que hace es pedir
``muestrame todo lo que tienes'' (esto se llama \textbf{service
discovery}). El servidor responde con la lista completa de
servicios, caracteristicas y descriptores.

\subsubsection{Servicio (Service)}

Un servicio es un \textbf{grupo logico de datos relacionados}.
En nuestro caso, creamos un servicio llamado ``IMU'' que contiene
las caracteristicas del acelerometro y el giroscopio. Si tuvieramos
tambien un sensor de temperatura, podriamos crear otro servicio
``Environment'' con una caracteristica de temperatura.

Cada servicio se identifica por un \textbf{UUID} (Universally Unique
Identifier --- ``Identificador Unico Universal''). Es un numero de
128 bits (16 bytes) que sirve como nombre unico.

Existen dos tipos de UUIDs:

\begin{itemize}
    \item \textbf{UUIDs estandar (16 bits)}: Definidos oficialmente
          por el Bluetooth SIG (la organizacion que gestiona el
          estandar Bluetooth). Son UUIDs cortos que se expanden a
          128 bits con una base fija. Ejemplo:
          \texttt{0x180F} = servicio de bateria,
          \texttt{0x180D} = servicio de ritmo cardiaco. Solo los
          puede asignar el Bluetooth SIG.

    \item \textbf{UUIDs custom (128 bits)}: Los creas tu. Como hay
          $2^{128}$ combinaciones posibles (un numero de 39 digitos),
          la probabilidad de colision es practicamente cero. Se
          escriben en formato
          \texttt{XXXXXXXX-XXXX-XXXX-XXXX-XXXXXXXXXXXX}.
          Puedes generar uno con el comando \texttt{uuidgen} en la
          terminal de tu Mac.
\end{itemize}

Nosotros usamos UUIDs custom porque no existe un servicio estandar
para ``datos raw de IMU''. El UUID que usamos es
\texttt{19B10000-E8F2-537E-4F6C-D104768A1214}. Lo importante no es
el numero en si, sino que sea \textbf{el mismo en Arduino y en
Python}.

\subsubsection{Caracteristica (Characteristic)}

Es la \textbf{unidad minima de dato}. Cada caracteristica contiene:

\begin{enumerate}
    \item \textbf{UUID}: Para identificarla dentro del servicio.
          Usamos \texttt{19B10001-...} para el acelerometro y
          \texttt{19B10002-...} para el giroscopio.

    \item \textbf{Valor}: Los bytes de datos propiamente dichos. En
          nuestro caso, 12 bytes que representan tres numeros float
          (ax, ay, az o gx, gy, gz).

    \item \textbf{Propiedades}: Flags (banderas) que definen que
          operaciones estan permitidas:

    \begin{itemize}
        \item \texttt{READ}: El central puede preguntar
              ``¿cual es tu valor ahora?'' y la caracteristica
              responde con sus bytes. Es como mirar un termometro:
              tu miras y lees el numero.

        \item \texttt{WRITE}: El central puede cambiar el valor. Se
              usa para enviar comandos al periferico (ej: ``cambia
              la frecuencia de muestreo a 50 Hz''). Nosotros no lo
              usamos.

        \item \texttt{NOTIFY}: La caracteristica envia su valor
              \textbf{automaticamente} cada vez que cambia, sin que
              el central lo pida. Es como suscribirte a un canal de
              noticias: te llegan las actualizaciones solas.
              \textbf{Esto es lo que usamos para streaming.}

        \item \texttt{INDICATE}: Similar a NOTIFY pero con
              confirmacion. El central responde ``recibido'' cada vez.
              Mas lento pero garantiza entrega. No lo usamos porque
              la velocidad es mas importante.
    \end{itemize}
\end{enumerate}

\subsubsection{Descriptor (Descriptor)}

Es metadata \emph{sobre} la caracteristica. El descriptor mas
importante es el \textbf{CCCD} (Client Characteristic Configuration
Descriptor, UUID \texttt{0x2902}).

El CCCD es como un ``interruptor'' de notificaciones. Es un valor
de 2 bytes que el central escribe para decir:

\begin{itemize}
    \item \texttt{0x0000}: ``No quiero notificaciones''
          (apagado).
    \item \texttt{0x0001}: ``Quiero recibir notificaciones''
          (encendido).
    \item \texttt{0x0002}: ``Quiero recibir indicaciones''
          (con confirmacion).
\end{itemize}

Cuando tu llamas a \texttt{client.start\_notify()} en Python, bleak
escribe \texttt{0x0001} en el CCCD de esa caracteristica. Cuando
llamas a \texttt{client.stop\_notify()}, escribe \texttt{0x0000}.

\textbf{¿Por que existe el CCCD?} Porque seria un desperdicio de
energia que el periferico enviara datos constantemente si nadie los
quiere. El CCCD permite al central decir ``estoy listo para
recibir'' antes de que el periferico empiece a gastar energia
enviando.

\subsection{Conexion BLE paso a paso}

Cuando tu script Python se conecta a la XIAO, esto es lo que ocurre
por debajo:

\begin{enumerate}
    \item \textbf{Scan}: Tu Mac enciende su receptor Bluetooth y
          escucha paquetes de advertising. Encuentra uno que dice
          ``Me llamo XIAO-IMU''.

    \item \textbf{Connection Request}: Tu Mac envia un paquete
          \texttt{CONNECT\_REQ} a la XIAO con los parametros de
          conexion deseados (intervalo, timeout, etc.).

    \item \textbf{Connection Established}: La XIAO acepta. Ahora
          ambos dejan de usar advertising y empiezan a comunicarse
          en slots de tiempo reservados (esto se llama
          ``connection events'').

    \item \textbf{MTU Exchange}: Negocian el tamaño maximo de
          paquete. MTU = Maximum Transmission Unit. Por defecto
          son 23 bytes (20 de payload + 3 de cabecera). Se puede
          negociar hasta 512 bytes.

    \item \textbf{Service Discovery}: Tu Mac le pide a la XIAO
          ``dime todos tus servicios y caracteristicas''. La XIAO
          responde con la tabla GATT completa. Esto solo ocurre
          una vez, al inicio.

    \item \textbf{Enable Notifications}: Tu Mac escribe
          \texttt{0x0001} en el CCCD de las caracteristicas que le
          interesan.

    \item \textbf{Data Streaming}: A partir de ahora, cada vez que
          la XIAO actualiza el valor de una caracteristica, envia
          un paquete de notificacion automaticamente.
\end{enumerate}

\subsection{Connection interval (intervalo de conexion)}

Una vez conectados, el central y el periferico no estan comunicandose
todo el tiempo. Eso consumiria demasiada energia. En su lugar,
``despiertan'' el radio cada cierto intervalo, intercambian datos
rapidamente, y vuelven a ``dormir''.

Este intervalo se llama \textbf{connection interval} y se mide en
unidades de 1.25 milisegundos. En nuestro codigo Arduino ponemos:

\begin{lstlisting}[style=arduino]
BLE.setConnectionInterval(8, 16);
\end{lstlisting}

Esto significa:
\begin{itemize}
    \item Minimo: $8 \times 1.25 = 10$ ms
    \item Maximo: $16 \times 1.25 = 20$ ms
\end{itemize}

Es decir, cada 10--20 ms hay una oportunidad de enviar datos. Como
nuestro sensor muestrea a $\sim$100 Hz (un dato cada 10 ms), un
intervalo de 10 ms nos permite enviar cada muestra en tiempo real.

\textbf{Nota:} Esto es una \emph{solicitud}, no una orden. El
sistema operativo de tu Mac puede negociar un intervalo diferente.
macOS, por ejemplo, raramente baja de 15 ms.


% ============================================================
\section{Numeros flotantes y bytes: como viajan los datos}
% ============================================================

\subsection{El problema: BLE solo envía bytes}

BLE no sabe que es un ``float'' ni un ``entero''. Solo sabe enviar
\textbf{bytes} (numeros enteros de 0 a 255). Si queremos enviar el
numero $1.0234$ (un float), necesitamos convertirlo a una secuencia
de bytes, enviarlo, y reconvertirlo en el otro lado.

\subsection{IEEE 754: como un float se convierte en 4 bytes}

Los ordenadores almacenan numeros decimales usando el estandar
\textbf{IEEE 754}. Un \texttt{float} de 32 bits (4 bytes) se
descompone asi:

\begin{verbatim}
  Bit:  31   30-23      22-0
        [S] [EEEEEEEE] [MMMMMMMMMMMMMMMMMMMMMMM]
         |   exponente   mantisa (parte fraccionaria)
         signo
\end{verbatim}

\begin{itemize}
    \item \textbf{S} (1 bit): 0 = positivo, 1 = negativo.
    \item \textbf{E} (8 bits): Exponente. Determina el ``orden de
          magnitud''. Se almacena con un sesgo de 127 (si E = 130,
          el exponente real es $130 - 127 = 3$).
    \item \textbf{M} (23 bits): Mantisa. La parte fraccionaria.
          Hay un ``1.'' implicito delante, asi que 23 bits dan
          24 bits efectivos de precision (unas 7 cifras decimales).
\end{itemize}

\textbf{Ejemplo con el numero 1.0:}
\begin{itemize}
    \item Signo: 0 (positivo)
    \item Exponente: 127 (es decir, $2^0 = 1$). En binario: 01111111
    \item Mantisa: 000...0 (porque $1.0 = 1.0 \times 2^0$, y la
          parte fraccionaria es cero)
\end{itemize}
Resultado en binario: \texttt{0 01111111 00000000000000000000000}

En hexadecimal: \texttt{0x3F800000}

En bytes: \texttt{[0x3F, 0x80, 0x00, 0x00]}

\subsection{Little-endian y Big-endian}

Los 4 bytes de un float se pueden ordenar de dos formas:

\begin{itemize}
    \item \textbf{Big-endian}: El byte mas significativo va primero.
          El numero \texttt{0x3F800000} se almacena como
          \texttt{[3F, 80, 00, 00]}.
    \item \textbf{Little-endian}: El byte menos significativo va
          primero. \texttt{[00, 00, 80, 3F]}.
\end{itemize}

El procesador ARM Cortex-M4 (que esta dentro del nRF52840) usa
\textbf{little-endian}. Los procesadores Intel/AMD y Apple Silicon
de tu Mac tambien. Esto es conveniente: los bytes llegan en el
mismo orden en que fueron escritos, sin necesidad de reordenar.

En Python, indicamos el orden de bytes en el formato de
\texttt{struct.unpack}:
\begin{itemize}
    \item \texttt{'<'} = little-endian
    \item \texttt{'>'} = big-endian
    \item \texttt{'f'} = un float de 32 bits (4 bytes)
    \item \texttt{'<fff'} = tres floats little-endian = 12 bytes
\end{itemize}

\subsection{memcpy: copiar bytes crudos en Arduino}

En el codigo Arduino, usamos \texttt{memcpy} para ``empaquetar''
los floats en un array de bytes:

\begin{lstlisting}[style=arduino]
float ax = 1.0;           // 4 bytes en memoria
uint8_t buffer[12];        // 12 bytes vacios
memcpy(&buffer[0], &ax, 4); // copia los 4 bytes de ax al buffer
\end{lstlisting}

\texttt{memcpy} significa ``memory copy'' (copia de memoria). Sus
parametros son:
\begin{enumerate}
    \item \textbf{destino}: \texttt{\&buffer[0]} --- la direccion de
          memoria donde queremos copiar (posicion 0 del array).
    \item \textbf{origen}: \texttt{\&ax} --- la direccion de memoria
          donde esta el float. El \texttt{\&} significa ``la direccion
          de''.
    \item \textbf{cantidad}: \texttt{4} --- cuantos bytes copiar (un
          float = 4 bytes).
\end{enumerate}

\texttt{memcpy} no sabe ni le importa que hay en esos bytes. Solo
copia bytes de un sitio a otro. Esto es lo que necesitamos: tomar
los 4 bytes que representan el float en memoria y ponerlos en
nuestro buffer de envio BLE.


% ============================================================
\section{Las librerias: quien hizo que y como funciona por dentro}
% ============================================================

\subsection{Lado Arduino (firmware de la XIAO)}

\subsubsection{ArduinoBLE}

\begin{itemize}
    \item \textbf{Proveedor}: Arduino (empresa/fundacion).
    \item \textbf{Repositorio}:
          \url{https://github.com/arduino-libraries/ArduinoBLE}
    \item \textbf{Que hace}: Proporciona toda la funcionalidad BLE
          en Arduino. Te permite crear servicios, caracteristicas,
          hacer advertising, gestionar conexiones, etc.
    \item \textbf{Como funciona por dentro}: ArduinoBLE es una capa
          de abstraccion sobre el \textbf{SoftDevice} de Nordic
          Semiconductor. El SoftDevice es un firmware propietario
          (un blob binario, no puedes ver su codigo fuente) que
          Nordic instala en parte de la memoria flash del nRF52840.
          Este firmware controla directamente el hardware de radio
          (el transceptor de 2.4 GHz, el sintetizador de frecuencia,
          etc.). ArduinoBLE llama a las funciones del SoftDevice a
          traves de la capa mbed-OS.

          La cadena completa es:

\begin{verbatim}
  Tu codigo Arduino
    -> ArduinoBLE (API amigable)
      -> mbed-OS cordio (capa de abstraccion)
        -> SoftDevice de Nordic (firmware del radio)
          -> Hardware de radio del nRF52840
\end{verbatim}
\end{itemize}

\subsubsection{Seeed\_Arduino\_LSM6DS3}

\begin{itemize}
    \item \textbf{Proveedor}: Seeed Studio (la empresa que fabrica
          la XIAO).
    \item \textbf{Repositorio}:
          \url{https://github.com/Seeed-Studio/Seeed_Arduino_LSM6DS3}
    \item \textbf{Que hace}: Permite leer datos del sensor
          LSM6DS3TR-C de forma sencilla con funciones como
          \texttt{readFloatAccelX()}.
    \item \textbf{Como funciona por dentro}: Es un wrapper (envoltorio)
          sobre el driver oficial de STMicroelectronics. Internamente
          usa la libreria Wire para comunicarse por I$^2$C:

\begin{verbatim}
  myIMU.readFloatAccelX()
    -> LSM6DS3::readRegisterInt16(OUTX_L_XL)
      -> Wire.beginTransmission(0x6A)
      -> Wire.write(0x28)     // registro a leer
      -> Wire.endTransmission()
      -> Wire.requestFrom(0x6A, 2)  // pedir 2 bytes
      -> byte_low = Wire.read()
      -> byte_high = Wire.read()
    -> raw_value = (byte_high << 8) | byte_low
    -> return raw_value * 0.000061  // convertir a g
\end{verbatim}
\end{itemize}

\subsubsection{Wire}

\begin{itemize}
    \item \textbf{Proveedor}: Arduino (viene incluida con Arduino
          IDE, es una ``core library'').
    \item \textbf{Que hace}: Implementa el protocolo I$^2$C.
    \item \textbf{Como funciona por dentro}: Configura los registros
          de hardware del modulo I$^2$C del nRF52840 (llamado TWIM
          --- Two-Wire Interface Master). Gestiona el timing de SCL,
          los ACK/NACK, y las interrupciones.
\end{itemize}

\subsection{Lado Python (tu ordenador)}

\subsubsection{bleak}

\begin{itemize}
    \item \textbf{Proveedor}: Henrik Blidh (desarrollador
          independiente, proyecto open source).
    \item \textbf{Repositorio}: \url{https://github.com/hbldh/bleak}
    \item \textbf{PyPI}: \url{https://pypi.org/project/bleak/}
    \item \textbf{Que hace}: Permite a Python comunicarse con
          dispositivos BLE. Es la libreria estandar de facto para
          BLE en Python.
    \item \textbf{Como funciona por dentro}: bleak no controla el
          Bluetooth directamente. Se apoya en las APIs del sistema
          operativo:

    \begin{itemize}
        \item \textbf{macOS}: Usa \texttt{CoreBluetooth}, el framework
              de Apple para BLE. Se comunica con el a traves de
              \texttt{pyobjc} (un bridge que permite llamar a APIs
              de Objective-C desde Python). Cuando llamas a
              \texttt{BleakScanner.discover()}, bleak le dice a
              CoreBluetooth ``escanea dispositivos BLE'', y
              CoreBluetooth le pide al chip Bluetooth del Mac que
              encienda el receptor y escuche paquetes de advertising.

        \item \textbf{Linux}: Usa \texttt{BlueZ} (el stack Bluetooth
              del kernel Linux) a traves de D-Bus (un sistema de
              comunicacion entre procesos).

        \item \textbf{Windows}: Usa \texttt{WinRT} (Windows Runtime
              APIs) para acceder al Bluetooth.
    \end{itemize}

    La cadena en macOS es:

\begin{verbatim}
  Tu codigo Python
    -> bleak (API asíncrona en Python)
      -> pyobjc (bridge Python-Objective-C)
        -> CoreBluetooth (framework de Apple)
          -> Driver Bluetooth de macOS
            -> Chip Bluetooth del Mac
\end{verbatim}

    \textbf{Toda la API de bleak es asincrona}. Esto significa que
    usa \texttt{asyncio} y que todas las operaciones importantes
    usan \texttt{async/await}. ¿Por que? Porque las operaciones BLE
    (escanear, conectar, esperar notificaciones) pueden tardar un
    tiempo indefinido, y no queremos que el programa se ``congele''
    mientras espera.
\end{itemize}

\subsubsection{struct (biblioteca estandar de Python)}

\begin{itemize}
    \item \textbf{Proveedor}: Python Software Foundation (viene
          incluido con Python, no hay que instalar nada).
    \item \textbf{Que hace}: Convierte entre valores de Python (int,
          float, etc.) y secuencias de bytes en formato C.
    \item \textbf{Funcion que usamos}: \texttt{struct.unpack(formato,
          bytes)}

          Toma una secuencia de bytes y la descompone en valores de
          Python segun el formato especificado. El formato
          \texttt{'<fff'} significa:
    \begin{itemize}
        \item \texttt{<}: Los bytes estan en orden little-endian.
        \item \texttt{f}: Un float de 32 bits (4 bytes).
        \item Tres \texttt{f}: Tres floats = 12 bytes total.
    \end{itemize}

          Devuelve una \textbf{tupla} (tipo de dato inmutable de
          Python) con los tres valores:
          \texttt{(1.0234, -0.0012, 0.9876)}
\end{itemize}

\subsubsection{asyncio (biblioteca estandar de Python)}

\begin{itemize}
    \item \textbf{Proveedor}: Python Software Foundation (incluido
          con Python 3.4+).
    \item \textbf{Que hace}: Permite escribir codigo asíncrono
          (que puede ``esperar'' sin bloquear el programa).
    \item \textbf{Conceptos clave}:

    \begin{itemize}
        \item \textbf{Event loop} (bucle de eventos): Es el
              ``corazon'' de asyncio. Es un bucle infinito que
              revisa constantemente: ``¿hay algun evento nuevo?
              ¿alguna operacion termino? ¿llego algun dato?''. Cuando
              detecta algo, ejecuta el codigo correspondiente.

              Analogia: imagina un camarero en un restaurante. En
              vez de quedarse parado esperando a que la cocina
              termine un plato, va a tomar nota de otras mesas.
              Cuando la cocina avisa que un plato esta listo, el
              camarero lo lleva. El event loop es ese camarero.

        \item \textbf{Coroutine}: Una funcion definida con
              \texttt{async def}. A diferencia de una funcion
              normal, una coroutine puede ``pausarse'' (con
              \texttt{await}) y reanudarse mas tarde. Cuando se
              pausa, el event loop puede ejecutar otras coroutines.

        \item \textbf{await}: Pausa la coroutine actual hasta que
              la operacion termine. Durante la pausa, el event loop
              puede hacer otras cosas (como procesar notificaciones
              BLE entrantes).

        \item \textbf{asyncio.run(coroutine)}: Crea un event loop,
              ejecuta la coroutine, y limpia todo al terminar. Es
              el punto de entrada para cualquier programa asincrono.

        \item \textbf{asyncio.sleep(segundos)}: Pausa la coroutine
              durante N segundos \textbf{sin bloquear el event loop}.
              Si usaras \texttt{time.sleep()} en su lugar,
              \emph{bloquearias} el event loop y no se procesarian
              las notificaciones BLE mientras dure la pausa.
    \end{itemize}
\end{itemize}


% ============================================================
\part{El codigo explicado linea por linea}
% ============================================================


% ============================================================
\section{Codigo Arduino completo}
% ============================================================

\subsection{Linea 1: \texttt{\#include <ArduinoBLE.h>}}

\begin{lstlisting}[style=arduino]
#include <ArduinoBLE.h>
\end{lstlisting}

\texttt{\#include} es una directiva del preprocesador de C++. Antes
de compilar tu codigo, el preprocesador reemplaza esta linea por
\textbf{todo el contenido} del archivo \texttt{ArduinoBLE.h}. Es
como hacer copy-paste automatico de ese archivo.

Los angulos \texttt{< >} indican que el archivo esta en las carpetas
del sistema (librerias instaladas). Si usaramos comillas
\texttt{" "}, buscaria en la carpeta del proyecto.

Al incluir \texttt{ArduinoBLE.h}, traemos al scope (ambito) las
siguientes clases y objetos:

\begin{itemize}
    \item \texttt{BLE}: Un objeto \textbf{singleton} global
          (singleton = solo existe uno en todo el programa). Controla
          el radio BLE del chip. Con el inicializas BLE, pones nombre,
          empiezas advertising, etc.

    \item \texttt{BLEService}: Clase para crear servicios GATT.
          Un servicio es un contenedor de caracteristicas (como
          explicamos antes).

    \item \texttt{BLECharacteristic}: Clase para crear
          caracteristicas GATT. Cada caracteristica almacena un
          valor (array de bytes) y tiene propiedades (READ, NOTIFY,
          etc.).

    \item \texttt{BLEDevice}: Clase que representa un dispositivo
          BLE conectado. Cuando un central se conecta a tu XIAO,
          obtienes un objeto \texttt{BLEDevice} que representa a
          ese central.
\end{itemize}

\subsection{Linea 2: \texttt{\#include "LSM6DS3.h"}}

\begin{lstlisting}[style=arduino]
#include "LSM6DS3.h"
\end{lstlisting}

Incluye la libreria del sensor IMU. Trae la clase \texttt{LSM6DS3}
que abstrae toda la comunicacion I$^2$C con el chip sensor. Sin
esta libreria, tendriamos que escribir manualmente los comandos
I$^2$C para leer cada registro (decenas de lineas de codigo).

Las comillas \texttt{" "} en vez de \texttt{< >} son una convencion:
indican que es una libreria externa (no del core de Arduino),
aunque en la practica ambas sintaxis funcionan igual.

\subsection{Linea 3: \texttt{\#include "Wire.h"}}

\begin{lstlisting}[style=arduino]
#include "Wire.h"
\end{lstlisting}

Incluye la libreria Wire que implementa I$^2$C. Nosotros no la
usamos directamente en nuestro codigo, pero la libreria LSM6DS3
la necesita internamente. Sin este include, el compilador daria
un error diciendo que no encuentra las funciones de Wire.

\subsection{Linea 4: Crear el objeto IMU}

\begin{lstlisting}[style=arduino]
LSM6DS3 myIMU(I2C_MODE, 0x6A);
\end{lstlisting}

Crea una \textbf{instancia} (un objeto concreto) de la clase
\texttt{LSM6DS3}. Le damos el nombre \texttt{myIMU}. Es como decir
``crea un sensor de tipo LSM6DS3 y llamalo myIMU''.

El \textbf{constructor} (la funcion que se ejecuta al crear el
objeto) recibe dos parametros:

\begin{enumerate}
    \item \texttt{I2C\_MODE}: Una constante definida en la libreria.
          Le dice al objeto que se comunique por I$^2$C (la
          alternativa seria \texttt{SPI\_MODE}, otro protocolo de
          comunicacion, pero la XIAO usa I$^2$C).

    \item \texttt{0x6A}: La direccion I$^2$C del sensor. El
          \texttt{0x} indica que es un numero en hexadecimal.
          \texttt{6A} en hexadecimal = 106 en decimal. Esta direccion
          viene predeterminada de fabrica en el chip.
\end{enumerate}

\textbf{Importante}: Crear el objeto no inicializa el sensor. Solo
guarda los parametros en memoria. La inicializacion real ocurre
cuando llamas a \texttt{myIMU.begin()}.

Internamente, el constructor hace algo como:
\begin{lstlisting}[style=arduino]
// Pseudocodigo interno del constructor
this->commMode = I2C_MODE;
this->I2CAddress = 0x6A;
this->accelRange = 2;     // default +/- 2g
this->gyroRange = 245;    // default +/- 245 dps
this->accelODR = 416;     // default 416 Hz
this->gyroODR = 416;      // default 416 Hz
\end{lstlisting}

\subsection{Lineas 5--7: Crear servicio y caracteristicas BLE}

\begin{lstlisting}[style=arduino]
BLEService imuService("19B10000-E8F2-537E-4F6C-D104768A1214");
\end{lstlisting}

Crea un servicio GATT. El parametro es el \textbf{UUID del servicio}
como string. Este UUID es un identificador unico que permite a los
clientes (tu Python) encontrar este servicio especifico entre todos
los servicios del dispositivo.

\begin{lstlisting}[style=arduino]
BLECharacteristic accelChar(
    "19B10001-E8F2-537E-4F6C-D104768A1214",
    BLERead | BLENotify,
    12
);
\end{lstlisting}

Crea una caracteristica GATT. Tres parametros:

\begin{enumerate}
    \item \textbf{UUID} (\texttt{"19B10001-..."}): Fijate que cambia
          el \texttt{0001} respecto al servicio (\texttt{0000}). Esto
          es una convencion para mantener los UUIDs relacionados.

    \item \textbf{Propiedades} (\texttt{BLERead | BLENotify}):
          \texttt{BLERead} y \texttt{BLENotify} son constantes
          numericas (por ejemplo, \texttt{BLERead = 0x02} y
          \texttt{BLENotify = 0x10}). El operador \texttt{|}
          (OR de bits) las combina: $0x02 | 0x10 = 0x12$. Esto
          crea un unico byte donde los bits de READ y NOTIFY estan
          activados.

    \item \textbf{Tamaño} (\texttt{12}): Tamaño maximo del valor en
          bytes. Necesitamos 12 porque enviamos 3 floats:
          $3 \times 4 = 12$ bytes.
\end{enumerate}

La caracteristica del giroscopio es identica pero con UUID
\texttt{0002}:

\begin{lstlisting}[style=arduino]
BLECharacteristic gyroChar(
    "19B10002-E8F2-537E-4F6C-D104768A1214",
    BLERead | BLENotify,
    12
);
\end{lstlisting}

\subsection{Funcion setup(): inicializacion}

\begin{lstlisting}[style=arduino]
void setup() {
\end{lstlisting}

\texttt{setup()} es una funcion especial de Arduino que se ejecuta
\textbf{una sola vez} al encender la placa. Aqui va toda la
inicializacion. \texttt{void} significa que la funcion no devuelve
ningun valor.

\begin{lstlisting}[style=arduino]
  Serial.begin(115200);
\end{lstlisting}

Inicializa la comunicacion serie por USB. \texttt{Serial} es un
objeto global que representa el puerto USB de la XIAO. Cuando la
XIAO esta conectada por USB al ordenador, puedes ver mensajes de
debug en el ``Monitor Serie'' del Arduino IDE.

\texttt{115200} es la velocidad en \textbf{baudios} (bits por
segundo). Es una velocidad estandar rapida. Tanto el emisor (XIAO)
como el receptor (tu ordenador) deben usar la misma velocidad.

\begin{lstlisting}[style=arduino]
  if (myIMU.begin() != 0) {
    Serial.println("IMU error");
    while (1);
  }
\end{lstlisting}

\texttt{myIMU.begin()} inicializa el sensor. Internamente hace:

\begin{enumerate}
    \item Llama a \texttt{Wire.begin()} para inicializar el bus
          I$^2$C (configura los pines SDA/SCL, activa el modulo
          de hardware I$^2$C del nRF52840).

    \item Lee el registro \texttt{WHO\_AM\_I} (direccion
          \texttt{0x0F}) del sensor. Este registro siempre contiene
          el valor \texttt{0x6A}. Es como preguntarle al sensor
          ``¿quien eres?'' para confirmar que esta ahi y funciona.

    \item Escribe en los registros de configuracion del sensor
          para establecer los rangos, frecuencias de muestreo, y
          filtros por defecto.

    \item Devuelve \texttt{0} si todo fue bien, o un numero distinto
          de cero si algo fallo (por ejemplo, si el sensor no
          respondio al \texttt{WHO\_AM\_I}).
\end{enumerate}

Si falla (\texttt{!= 0}), imprime un mensaje de error y entra en
\texttt{while(1)}: un bucle infinito que detiene el programa. Esto
es un patron comun en embedded: si algo critico falla en la
inicializacion, es mejor parar que seguir con un sistema roto.

\begin{lstlisting}[style=arduino]
  if (!BLE.begin()) {
    Serial.println("BLE failed!");
    while (1);
  }
\end{lstlisting}

\texttt{BLE.begin()} inicializa el stack BLE. Internamente:

\begin{enumerate}
    \item Inicializa el SoftDevice de Nordic (el firmware
          propietario que controla el radio de 2.4 GHz).
    \item Reserva memoria RAM para las estructuras BLE (tabla GATT,
          buffers de conexion, colas de paquetes).
    \item Configura las interrupciones de hardware para que el
          radio pueda avisar al procesador cuando llega un paquete.
    \item Devuelve \texttt{true} si todo fue bien, \texttt{false}
          si fallo.
\end{enumerate}

El \texttt{!} (NOT) niega el valor: si \texttt{begin()} devuelve
\texttt{false} (fallo), \texttt{!false = true}, y entramos al
bloque de error.

\begin{lstlisting}[style=arduino]
  BLE.setLocalName("XIAO-IMU");
\end{lstlisting}

Establece el nombre que aparecera en los paquetes de advertising.
Cuando hagas un scan BLE desde tu Mac o telefono, veras ``XIAO-IMU''
en la lista. Este nombre se incluye en el paquete de advertising
como un campo de tipo ``Complete Local Name''. Limite: 29 caracteres.

\begin{lstlisting}[style=arduino]
  BLE.setAdvertisedService(imuService);
\end{lstlisting}

Le dice al stack BLE que incluya el UUID de nuestro servicio IMU en
los paquetes de advertising. Esto permite que los centrales puedan
filtrar por servicio antes de conectarse. Sin esta linea, el
dispositivo se anunciaria sin decir que servicios tiene, y el central
tendria que conectarse ``a ciegas'' para descubrirlo.

\begin{lstlisting}[style=arduino]
  imuService.addCharacteristic(accelChar);
  imuService.addCharacteristic(gyroChar);
\end{lstlisting}

Asocia las caracteristicas al servicio. Construye la jerarquia GATT.
Internamente, el servicio mantiene una lista enlazada de
caracteristicas. Al añadir una, se agrega al final de esa lista.

\begin{lstlisting}[style=arduino]
  BLE.addService(imuService);
\end{lstlisting}

Registra el servicio completo (con todas sus caracteristicas) en la
tabla GATT del stack BLE. A partir de ahora, el servicio ``existe''
y sera visible para cualquier central que se conecte y haga service
discovery.

\begin{lstlisting}[style=arduino]
  BLE.setConnectionInterval(8, 16);
\end{lstlisting}

Ya lo explicamos en la seccion de connection interval. Pedimos un
intervalo de 10--20 ms.

\begin{lstlisting}[style=arduino]
  BLE.advertise();
\end{lstlisting}

Empieza a emitir paquetes de advertising. A partir de este momento,
la XIAO es visible para cualquier scanner BLE. Los paquetes se
emiten aproximadamente cada 100 ms hasta que alguien se conecte.

\subsection{Funcion loop(): el bucle principal}

\begin{lstlisting}[style=arduino]
void loop() {
  BLEDevice central = BLE.central();
\end{lstlisting}

\texttt{loop()} se ejecuta en bucle infinito despues de
\texttt{setup()}. Es el ``main loop'' del programa.

\texttt{BLE.central()} comprueba si hay un dispositivo central
conectado. Devuelve un objeto \texttt{BLEDevice}:
\begin{itemize}
    \item Si hay conexion: devuelve un objeto valido que contiene la
          direccion MAC del central.
    \item Si no hay conexion: devuelve un objeto ``vacio'' que
          evalua a \texttt{false} en un \texttt{if}.
\end{itemize}

\begin{lstlisting}[style=arduino]
  if (central) {
    Serial.println("Connected");
    while (central.connected()) {
\end{lstlisting}

Si hay un central conectado, entramos en un bucle \texttt{while}
que continua mientras la conexion este activa.
\texttt{central.connected()} devuelve \texttt{false} cuando:
\begin{itemize}
    \item El central se desconecta voluntariamente.
    \item Se pierde la conexion (la XIAO se aleja demasiado, hay
          interferencia, etc.).
    \item Pasan mas de 4 segundos sin comunicacion (supervision
          timeout).
\end{itemize}

\begin{lstlisting}[style=arduino]
      float ax = myIMU.readFloatAccelX();
      float ay = myIMU.readFloatAccelY();
      float az = myIMU.readFloatAccelZ();
      float gx = myIMU.readFloatGyroX();
      float gy = myIMU.readFloatGyroY();
      float gz = myIMU.readFloatGyroZ();
\end{lstlisting}

Lee los 6 valores del IMU. Cada llamada a \texttt{readFloat...()}:
\begin{enumerate}
    \item Envia un comando I$^2$C al sensor pidiendo 2 bytes del
          registro correspondiente.
    \item Recibe los 2 bytes (un entero de 16 bits con signo).
    \item Multiplica por el factor de escala para convertir a
          unidades fisicas.
    \item Devuelve un \texttt{float} (numero decimal de 32 bits).
\end{enumerate}

El acelerometro devuelve valores en \textbf{g} (1g $\approx$
9.81 m/s$^2$). En reposo: $a_x \approx 0, a_y \approx 0,
a_z \approx 1.0$.

El giroscopio devuelve valores en \textbf{grados/segundo}. En
reposo: $g_x \approx 0, g_y \approx 0, g_z \approx 0$.

\begin{lstlisting}[style=arduino]
      uint8_t accelBuf[12];
      memcpy(&accelBuf[0], &ax, 4);
      memcpy(&accelBuf[4], &ay, 4);
      memcpy(&accelBuf[8], &az, 4);
      accelChar.writeValue(accelBuf, 12);
\end{lstlisting}

\texttt{uint8\_t accelBuf[12]}: Declara un array de 12 bytes en la
pila (stack). \texttt{uint8\_t} = unsigned integer de 8 bits = un
byte (valor de 0 a 255).

Las tres lineas de \texttt{memcpy} copian los bytes crudos de cada
float al buffer. Despues de ejecutarse, el buffer contiene:

\begin{verbatim}
  Posicion:  0  1  2  3    4  5  6  7    8  9 10 11
  Contenido: [--- ax ---]  [--- ay ---]  [--- az ---]
             4 bytes       4 bytes       4 bytes
\end{verbatim}

\texttt{accelChar.writeValue(accelBuf, 12)}: Actualiza el valor de
la caracteristica GATT. Internamente:
\begin{enumerate}
    \item Copia los 12 bytes al buffer interno de la caracteristica.
    \item Si un central esta suscrito a notificaciones (escribio
          0x0001 en el CCCD), encola un paquete de notificacion
          BLE.
    \item El paquete se transmitira por radio en el siguiente
          connection event (el proximo ``despertar'' del intervalo
          de conexion).
\end{enumerate}

El bloque del giroscopio es identico:

\begin{lstlisting}[style=arduino]
      uint8_t gyroBuf[12];
      memcpy(&gyroBuf[0], &gx, 4);
      memcpy(&gyroBuf[4], &gy, 4);
      memcpy(&gyroBuf[8], &gz, 4);
      gyroChar.writeValue(gyroBuf, 12);
\end{lstlisting}

\begin{lstlisting}[style=arduino]
      delay(10);
\end{lstlisting}

Pausa la ejecucion 10 milisegundos. Esto limita la frecuencia de
muestreo a aproximadamente $1/0.010 = 100$ Hz.

\textbf{¿Por que ``aproximadamente''?} Porque las lecturas I$^2$C
(6 transacciones) y los envios BLE (2 \texttt{writeValue}) tambien
tardan tiempo (1--3 ms en total). Asi que el ciclo real es de
$\sim$12 ms, dando $\sim$83 Hz. Para la mayoria de aplicaciones
de IMU esto es suficiente.

\begin{lstlisting}[style=arduino]
    }
    Serial.println("Disconnected");
  }
}
\end{lstlisting}

Cuando el central se desconecta, sale del \texttt{while} e imprime
``Disconnected''. El \texttt{loop()} vuelve a empezar: llama de
nuevo a \texttt{BLE.central()}, que no encuentra conexion, asi que
no entra al \texttt{if}, y el loop se repite. El advertising se
reanuda automaticamente cuando no hay conexion activa, asi que la
XIAO vuelve a ser visible para nuevas conexiones.


% ============================================================
\section{Codigo Python completo}
% ============================================================

\subsection{Imports}

\begin{lstlisting}[style=python]
import asyncio
\end{lstlisting}

Importa el modulo \texttt{asyncio}. Ya explicamos a fondo que es
asyncio en la seccion de librerias. En resumen: nos permite
escribir codigo que ``espera'' (conexion BLE, datos, etc.) sin
congelar el programa.

\begin{lstlisting}[style=python]
import struct
\end{lstlisting}

Importa el modulo \texttt{struct}. Lo usamos para convertir los
bytes que llegan por BLE de vuelta a numeros float.

\begin{lstlisting}[style=python]
import time
\end{lstlisting}

Importa el modulo \texttt{time}. Lo usamos para obtener timestamps
(\texttt{time.time()} devuelve los segundos desde el 1 de enero de
1970, con precision de microsegundos).

\begin{lstlisting}[style=python]
from bleak import BleakClient, BleakScanner
\end{lstlisting}

Importa dos clases de bleak:

\begin{itemize}
    \item \texttt{BleakScanner}: Para escanear dispositivos BLE
          cercanos. Es como ``mirar que hay en el aire''.
    \item \texttt{BleakClient}: Para conectarse a un dispositivo
          especifico y comunicarse con el (leer/escribir
          caracteristicas, suscribirse a notificaciones, etc.).
\end{itemize}

\subsection{Constantes}

\begin{lstlisting}[style=python]
ACCEL_UUID = "19b10001-e8f2-537e-4f6c-d104768a1214"
GYRO_UUID  = "19b10002-e8f2-537e-4f6c-d104768a1214"
\end{lstlisting}

Los UUIDs de las caracteristicas GATT. \textbf{Deben coincidir
exactamente} con los del sketch Arduino. Estan en minusculas porque
bleak normaliza todos los UUIDs a minusculas internamente.

Estos strings son la ``direccion'' que Python usara para decirle a
bleak ``quiero los datos de esta caracteristica especifica''.

\subsection{Variables globales}

\begin{lstlisting}[style=python]
latest_accel = (0, 0, 0)
latest_gyro = (0, 0, 0)
\end{lstlisting}

Tuplas que almacenan la ultima lectura recibida. Las inicializamos
en $(0, 0, 0)$. Se actualizan desde los callbacks cada vez que
llega una notificacion.

\textbf{¿Que es una tupla?} Es un tipo de dato de Python similar a
una lista, pero inmutable (no se puede modificar una vez creada).
Se escribe con parentesis: \texttt{(1, 2, 3)}. Usamos tuplas porque
\texttt{struct.unpack} devuelve tuplas.

\textbf{¿Por que \texttt{global}?} En Python, si quieres modificar
una variable definida fuera de una funcion desde dentro de esa
funcion, debes declararla como \texttt{global}. Sin esa declaracion,
Python crearia una variable local con el mismo nombre (y la global
no se modificaria).

\subsection{Callbacks (funciones que se llaman automaticamente)}

\begin{lstlisting}[style=python]
def accel_handler(_, data: bytearray):
    global latest_accel
    latest_accel = struct.unpack('<fff', data)
    ax, ay, az = latest_accel
    print(f"Accel: x={ax:.4f}  y={ay:.4f}  z={az:.4f} g")
\end{lstlisting}

Esta funcion se ejecuta \textbf{automaticamente} cada vez que la
XIAO envia una notificacion de la caracteristica del acelerometro.
Nosotros nunca la llamamos directamente; la registramos con
\texttt{start\_notify} y bleak la llama cuando llegan datos.

Parametros:
\begin{itemize}
    \item \texttt{\_}: Es el objeto
          \texttt{BleakGATTCharacteristic} que origino la
          notificacion. Lo nombramos \texttt{\_} porque no lo
          necesitamos. En Python, \texttt{\_} es una convencion
          para ``variable que no me importa''.

    \item \texttt{data: bytearray}: Los bytes crudos que llegaron
          por BLE. \texttt{bytearray} es un tipo de Python: una
          secuencia mutable de bytes. Contendra exactamente 12
          bytes (los 3 floats que empaquetamos en Arduino).
          El \texttt{: bytearray} es un ``type hint'' (anotacion
          de tipo): no cambia el comportamiento del codigo, solo
          documenta que tipo de dato se espera.
\end{itemize}

Cuerpo:
\begin{itemize}
    \item \texttt{global latest\_accel}: Declara que vamos a
          modificar la variable global.

    \item \texttt{struct.unpack('<fff', data)}: Toma los 12 bytes
          de \texttt{data} y los desempaqueta como 3 floats
          little-endian. Devuelve una tupla, por ejemplo
          \texttt{(0.0234, -0.0012, 0.9876)}.

    \item \texttt{ax, ay, az = latest\_accel}: ``Destructuring'' o
          ``unpacking''. Asigna cada elemento de la tupla a una
          variable individual. Es equivalente a:
\begin{lstlisting}[style=python]
ax = latest_accel[0]
ay = latest_accel[1]
az = latest_accel[2]
\end{lstlisting}

    \item \texttt{print(f"...")}: ``f-string'' (formatted string
          literal). Las expresiones entre \texttt{\{\}} se evaluan
          y se insertan en el string. \texttt{:.4f} formatea el
          float con 4 decimales.
\end{itemize}

El handler del giroscopio es identico pero con \texttt{:.2f}
(2 decimales):

\begin{lstlisting}[style=python]
def gyro_handler(_, data: bytearray):
    global latest_gyro
    latest_gyro = struct.unpack('<fff', data)
    gx, gy, gz = latest_gyro
    print(f"Gyro:  x={gx:.2f}  y={gy:.2f}  z={gz:.2f} deg/s")
\end{lstlisting}

\textbf{¿Como llega la notificacion al callback?} La cadena es:
\begin{enumerate}
    \item La XIAO actualiza el valor de la caracteristica y el radio
          BLE envia un paquete de notificacion.
    \item El chip Bluetooth de tu Mac recibe el paquete.
    \item CoreBluetooth (macOS) procesa el paquete y notifica a
          bleak a traves de un ``delegate'' de Objective-C.
    \item bleak pone el evento en la cola del event loop de asyncio.
    \item El event loop ejecuta tu callback con los datos.
\end{enumerate}

Todo esto ocurre en milisegundos.

\subsection{Funcion principal}

\begin{lstlisting}[style=python]
async def main():
\end{lstlisting}

Define \texttt{main} como una \textbf{coroutine asincrona}.
\texttt{async def} es como \texttt{def} pero indica que la funcion
puede usar \texttt{await} y puede pausarse/reanudarse.

\begin{lstlisting}[style=python]
    print("Scanning...")
    device = await BleakScanner.find_device_by_name(
        "XIAO-IMU", timeout=10.0
    )
\end{lstlisting}

Escanea dispositivos BLE buscando uno con nombre ``XIAO-IMU''.

\texttt{BleakScanner.find\_device\_by\_name()} es un metodo
\textbf{de clase} (se llama en la clase, no en una instancia). Es
asincrono, por eso usamos \texttt{await}.

Internamente:
\begin{enumerate}
    \item Crea un objeto \texttt{BleakScanner}.
    \item Llama a CoreBluetooth para empezar a escanear.
    \item CoreBluetooth activa el receptor Bluetooth y empieza a
          escuchar paquetes de advertising.
    \item Por cada paquete recibido, compara el campo
          ``Local Name'' con \texttt{"XIAO-IMU"}.
    \item Si encuentra coincidencia, para el scan y devuelve un
          objeto \texttt{BLEDevice} con la informacion del
          dispositivo.
    \item Si pasan 10 segundos sin encontrarlo, devuelve
          \texttt{None}.
\end{enumerate}

El \texttt{await} pausa la coroutine \texttt{main()} mientras el
scan esta activo. El event loop queda libre para procesar otros
eventos (aunque en este caso no hay otros).

\begin{lstlisting}[style=python]
    if not device:
        print("XIAO-IMU not found!")
        return
\end{lstlisting}

Si \texttt{device} es \texttt{None} (no se encontro), muestra un
mensaje y termina la funcion. \texttt{return} en una coroutine
asincrona la finaliza y devuelve \texttt{None}.

\begin{lstlisting}[style=python]
    print(f"Found: {device.address}")
\end{lstlisting}

Imprime la direccion del dispositivo. En macOS, la ``direccion''
no es una MAC real (Apple la oculta por privacidad) sino un UUID
de CoreBluetooth. En Linux si veras la MAC real (algo como
\texttt{A4:CF:12:8B:3E:D2}).

\begin{lstlisting}[style=python]
    async with BleakClient(device, timeout=20.0) as client:
\end{lstlisting}

Esto hace varias cosas a la vez:

\begin{enumerate}
    \item \textbf{\texttt{BleakClient(device, timeout=20.0)}}:
          Crea un objeto cliente BLE configurado para conectarse
          al dispositivo encontrado. \texttt{timeout=20.0} es el
          tiempo maximo en segundos para establecer la conexion.

    \item \textbf{\texttt{async with ... as client}}: Es un
          ``context manager asincrono''. Es un patron de Python que
          garantiza limpieza de recursos. Equivale a:

\begin{lstlisting}[style=python]
client = BleakClient(device, timeout=20.0)
await client.connect()   # se conecta
try:
    # ... tu codigo ...
finally:
    await client.disconnect()  # SIEMPRE desconecta
\end{lstlisting}

          La ventaja del \texttt{async with} es que la desconexion
          ocurre \textbf{siempre}, incluso si hay un error o
          excepcion dentro del bloque. Asi no dejas conexiones BLE
          abiertas ``colgadas''.
\end{enumerate}

Al entrar al bloque, bleak:
\begin{enumerate}
    \item Pide a CoreBluetooth que establezca conexion con el
          periferico.
    \item Espera la confirmacion de conexion.
    \item Ejecuta \textbf{service discovery}: le pide al periferico
          su tabla GATT completa y la almacena en
          \texttt{client.services}.
\end{enumerate}

\begin{lstlisting}[style=python]
        print(f"Connected: {client.is_connected}")
\end{lstlisting}

\texttt{client.is\_connected} es una propiedad (no una funcion,
por eso no lleva parentesis) que devuelve \texttt{True} si hay una
conexion BLE activa.

\begin{lstlisting}[style=python]
        await client.start_notify(ACCEL_UUID, accel_handler)
        await client.start_notify(GYRO_UUID, gyro_handler)
\end{lstlisting}

\textbf{Esta es la linea mas importante de todo el programa Python.}

\texttt{start\_notify(uuid, callback)} hace:

\begin{enumerate}
    \item Busca la caracteristica por UUID en la tabla GATT que se
          obtuvo durante service discovery.
    \item Verifica que la caracteristica tenga la propiedad NOTIFY
          (si no, lanza una excepcion).
    \item Escribe \texttt{0x0001} (2 bytes, little-endian) en el
          \textbf{CCCD} (descriptor \texttt{0x2902}) de esa
          caracteristica. Esto le dice a la XIAO: ``activa las
          notificaciones para esta caracteristica''.
    \item Registra \texttt{callback} (nuestra funcion
          \texttt{accel\_handler}) en un diccionario interno de
          bleak, indexado por el ``handle'' (identificador numerico
          interno) de la caracteristica.
    \item A partir de ahora, cada vez que la XIAO haga
          \texttt{writeValue()} en esa caracteristica, el stack BLE
          enviara los datos y bleak llamara a nuestro callback.
\end{enumerate}

\begin{lstlisting}[style=python]
        print("Streaming... Ctrl+C to stop\n")
\end{lstlisting}

Mensaje informativo. \texttt{\textbackslash n} es un salto de linea
extra.

\begin{lstlisting}[style=python]
        try:
            while True:
                await asyncio.sleep(1.0)
        except KeyboardInterrupt:
            pass
\end{lstlisting}

Este bloque mantiene el programa vivo mientras recibe datos:

\begin{itemize}
    \item \texttt{while True}: Bucle infinito. Sin esto, la funcion
          terminaria inmediatamente y se cerraria la conexion BLE.

    \item \texttt{await asyncio.sleep(1.0)}: Pausa la coroutine 1
          segundo. \textbf{Importantisimo}: usamos
          \texttt{asyncio.sleep} (asincrono) y NO \texttt{time.sleep}
          (sincrono). La diferencia:

    \begin{itemize}
        \item \texttt{asyncio.sleep}: Pausa \emph{solo esta
              coroutine}. El event loop sigue activo y puede
              procesar notificaciones BLE. Los callbacks se ejecutan
              normalmente durante la pausa.
        \item \texttt{time.sleep}: Bloquea \emph{todo el hilo},
              incluyendo el event loop. Ninguna notificacion BLE se
              procesaria durante la pausa. Los datos se perderian.
    \end{itemize}

    \item \texttt{except KeyboardInterrupt}: Captura la excepcion
          que se genera cuando presionas Ctrl+C en la terminal.
          \texttt{pass} significa ``no hacer nada especial, solo
          salir del bucle limpiamente''. Sin este try/except, al
          pulsar Ctrl+C verias un traceback de error feo.
\end{itemize}

\begin{lstlisting}[style=python]
        await client.stop_notify(ACCEL_UUID)
        await client.stop_notify(GYRO_UUID)
\end{lstlisting}

Cancela las suscripciones. Internamente:
\begin{enumerate}
    \item Escribe \texttt{0x0000} en el CCCD de cada caracteristica
          (``apaga notificaciones'').
    \item Elimina los callbacks del diccionario interno de bleak.
\end{enumerate}

No es estrictamente necesario porque \texttt{disconnect()} tambien
limpia todo, pero es buena practica hacer un cierre ordenado.

\begin{lstlisting}[style=python]
    print("Disconnected.")
\end{lstlisting}

Al salir del \texttt{async with}, bleak llama automaticamente a
\texttt{client.disconnect()}, que:
\begin{enumerate}
    \item Envia un paquete de terminacion de conexion BLE.
    \item Libera los recursos del radio Bluetooth en el SO.
    \item La XIAO detecta la desconexion y sale del
          \texttt{while(central.connected())}.
\end{enumerate}

\subsection{Punto de entrada}

\begin{lstlisting}[style=python]
if __name__ == "__main__":
    asyncio.run(main())
\end{lstlisting}

\begin{itemize}
    \item \texttt{if \_\_name\_\_ == "\_\_main\_\_"}: Comprueba si
          este archivo se esta ejecutando directamente (no
          importado como modulo). \texttt{\_\_name\_\_} es una
          variable especial que Python establece automaticamente:
    \begin{itemize}
        \item Si ejecutas \texttt{python script.py}:
              \texttt{\_\_name\_\_ = "\_\_main\_\_"}
        \item Si haces \texttt{import script}:
              \texttt{\_\_name\_\_ = "script"}
    \end{itemize}

    \item \texttt{asyncio.run(main())}: Punto de entrada del
          programa asincrono. Hace tres cosas:
    \begin{enumerate}
        \item Crea un nuevo event loop.
        \item Ejecuta la coroutine \texttt{main()} hasta que
              termine.
        \item Cierra el event loop y limpia todos los recursos.
    \end{enumerate}

    \textbf{Nota para Jupyter Notebook}: En un notebook, ya hay un
    event loop ejecutandose (el del propio notebook). Llamar a
    \texttt{asyncio.run()} dentro de un event loop existente da
    error. En su lugar, usa directamente:
\begin{lstlisting}[style=python]
await main()
\end{lstlisting}
\end{itemize}


% ============================================================
\part{Apendices}
% ============================================================


% ============================================================
\section{Extension: guardar datos en CSV}
% ============================================================

Para alimentar tus notebooks de Madgwick con datos BLE:

\begin{lstlisting}[style=python]
import csv

# Abrir archivo CSV al inicio
csv_file = open('data/ble_capture.csv', 'w', newline='')
writer = csv.writer(csv_file)
writer.writerow([
    'timestamp',
    'accel_x', 'accel_y', 'accel_z',
    'gyro_x', 'gyro_y', 'gyro_z'
])

def accel_handler(_, data: bytearray):
    global latest_accel
    latest_accel = struct.unpack('<fff', data)

def gyro_handler(_, data: bytearray):
    global latest_gyro
    latest_gyro = struct.unpack('<fff', data)
    # Escribir al CSV cuando llega el gyro
    t = time.time()
    ax, ay, az = latest_accel
    gx, gy, gz = latest_gyro
    writer.writerow([t, ax, ay, az, gx, gy, gz])
\end{lstlisting}


% ============================================================
\section{Diagnostico y solucion de problemas}
% ============================================================

\subsection{El dispositivo no aparece en el scan}

\begin{itemize}
    \item Verifica que el sketch Arduino se cargo correctamente
          (abre el Monitor Serie y busca ``Advertising as XIAO-IMU'').
    \item Asegura que no hay otro dispositivo ya conectado a la
          XIAO (BLE solo permite una conexion simultanea por defecto).
    \item En macOS: Preferencias del Sistema > Bluetooth > verificar
          que esta activado.
    \item El nombre es case-sensitive: ``XIAO-IMU'' $\neq$
          ``xiao-imu''.
    \item Aumenta el timeout del scan a 15--20 segundos.
\end{itemize}

\subsection{Se conecta pero no llegan datos}

\begin{itemize}
    \item Verifica que los UUIDs en Python son \textbf{identicos} a
          los del Arduino (copia y pega, no escribas a mano).
    \item Verifica que \texttt{myIMU.begin()} devolvio 0 (Monitor
          Serie).
    \item Ejecuta un script que liste servicios y caracteristicas
          para ver que expone realmente el dispositivo.
\end{itemize}

\subsection{Los valores no tienen sentido}

\begin{itemize}
    \item Verifica que usas \texttt{'<fff'} (little-endian) y no
          \texttt{'>fff'} (big-endian).
    \item Verifica que \texttt{len(data) == 12}. Si es diferente,
          el empaquetado no coincide.
    \item En reposo, el acelerometro debe leer $\approx (0, 0, 1)$ g
          y el giroscopio $\approx (0, 0, 0)$ dps.
\end{itemize}


% ============================================================
\section{Referencias}
% ============================================================

\begin{itemize}
    \item Bleak documentation: \url{https://bleak.readthedocs.io/}
    \item Bleak GitHub: \url{https://github.com/hbldh/bleak}
    \item ArduinoBLE reference:
          \url{https://www.arduino.cc/reference/en/libraries/arduinoble/}
    \item Seeed XIAO BLE Sense wiki --- IMU:
          \url{https://wiki.seeedstudio.com/XIAO-BLE-Sense-IMU-Usage/}
    \item Seeed XIAO BLE Sense wiki --- Bluetooth:
          \url{https://wiki.seeedstudio.com/XIAO-BLE-Sense-Bluetooth-Usage/}
    \item LSM6DS3 library:
          \url{https://github.com/Seeed-Studio/Seeed_Arduino_LSM6DS3}
    \item Bluetooth GATT specification:
          \url{https://www.bluetooth.com/specifications/gatt/}
    \item IEEE 754 floating point:
          \url{https://en.wikipedia.org/wiki/IEEE_754}
    \item Python struct module:
          \url{https://docs.python.org/3/library/struct.html}
    \item Python asyncio:
          \url{https://docs.python.org/3/library/asyncio.html}
\end{itemize}

\end{document}
